% !TeX root = ../navier-stokes-manuscript.tex
% ============================================================
% 2.1 Critical Height at a Finite Endpoint
% ============================================================

\subsection{Critical Height at a Finite Endpoint}\label{sub:CriticalHeightAtAFiniteEndpoint}

Fix a rapidly decaying datum \(u_0\) on \(\R^3\), a viscosity \(\nu>0\), and the maximal classical solution \((u,p)\) they generate on \([0,T_*)\).
Suppose \(T_*<\infty\).
The field is smooth at every earlier time, so any concentration has to be followed through smaller scales as \(t\) approaches that endpoint.

Kinetic energy stays bounded along the classical interval, but a fixed amount of energy is not equally visible at every scale.
The Navier--Stokes scaling shows the difference.
For every \(\lambda>0\), the rescaled fields
\[
  u_\lambda(x,t)
  :=
  \lambda u(\lambda x,\lambda^2t),
  \qquad
  p_\lambda(x,t)
  :=
  \lambda^2p(\lambda x,\lambda^2t)
\]
solve the same equation with the same viscosity on the rescaled time interval.
Length has been divided by \(\lambda\), time by \(\lambda^2\), and velocity has been multiplied by \(\lambda\).
This is the parabolic relation already present in the equation: one time derivative and the two spatial derivatives in viscosity acquire the same factor under rescaling.

Let \(\Lambda=(-\Delta)^{1/2}\), and measure a spatial Sobolev height by
\[
  \norm{v}_{\dot H^s}^2
  :=
  \int_{\R^3}
  |\xi|^{2s}|\widehat v(\xi)|^2
  \,d\xi.
\]
\begin{proposition}[Critical Sobolev scaling]\label[proposition]{prop:CriticalSobolevScaling}
For every \(s\in\R\) for which the homogeneous norm is finite and every \(\lambda>0\),
\[
  \norm{u_\lambda(t)}_{\dot H^s}^2
  =
  \lambda^{2s-1}
  \norm{u(\lambda^2t)}_{\dot H^s}^2.
\]
In particular, \(s=\tfrac12\) is the unique scale-invariant Sobolev height.
\end{proposition}
\begin{proof}
The Fourier transform of the rescaled velocity is
\[
  \widehat{u_\lambda}(\xi,t)
  =
  \lambda^{-2}
  \widehat u(\xi/\lambda,\lambda^2t).
\]
Putting this into the Sobolev norm and changing variables from \(\xi\) to \(\eta=\xi/\lambda\) gives
\begin{align*}
  \norm{u_\lambda(t)}_{\dot H^s}^2
  &={}
  \int_{\R^3}
  |\xi|^{2s}\lambda^{-4}
  \left|\widehat u(\xi/\lambda,\lambda^2t)\right|^2
  \,d\xi
  \\
  &={}
  \lambda^{2s-1}
  \int_{\R^3}
  |\eta|^{2s}
  \left|\widehat u(\eta,\lambda^2t)\right|^2
  \,d\eta.
\end{align*}
The exponent \(2s-1\) vanishes exactly at \(s=\tfrac12\), which proves the proposition.
\end{proof}
Every Sobolev height tilts under the compression except \(s=1/2\).
At that height the power of \(\lambda\) vanishes, so define
\[
  H(t)
  :=
  \frac12\norm{u(t)}_{\dot H^{1/2}}^2
  =
  \frac12\norm{\Lambda^{1/2}u(t)}_2^2.
\]
This is the critical height: exact rescaling can move a smooth configuration to a finer length and a shorter time without changing the value of \(H\).

The kinetic-energy scale moves in the other direction.
Putting \(s=0\) in the same calculation gives
\[
  \norm{u_\lambda(t)}_2^2
  =
  \lambda^{-1}
  \norm{u(\lambda^2t)}_2^2.
\]
A finer rescaled structure can carry less \(L^2\) energy while standing at the same critical height.
That is the opening left by the ordinary energy estimate: it can keep the total energy finite while the scale at which the fluid is changing continues to fall.

The finite parabolic schedule in \cref{sub:HistoricalEstimates} shows that progressively smaller scale durations can fit before one finite time.
At the height just derived, each finer copy can keep the same value of \(H\) while its \(L^2\) contribution shrinks.
Scaling alone does not make the actual fluid pass through those copies.
The equation would have to generate each finer state from the state before it while transport, pressure, incompressibility, and viscosity all continue to act.

Critical height fixes the scale.
The finite endpoint asks how fast the field can move through an unbounded sequence of regular states, which is a question about one time response and the response immediately above it.
